\centerline{\bf \Huge НОД и НОК. Алгоритм Евклида}
\title{НОД и НОК. Алгоритм Евклида.}

\begin{flushright}
    \scriptsize Я взял два простых числа. Два других простых числа. \\Перемножил. Получилось одно и то же!\\
    Райгородский. Анекдот про ферматистов.
\end{flushright}
\problem{0}Доказать следующие свойства НОД и НОК:

$\bullet \ (a, b)[a, b]=ab$

$\bullet \ d(m, \ n) = (dm, \ dn)$

$\bullet \  \text{Если}\ (m, \ n) = d, \ \text{то} \ (\dfrac{m}{d}, \ \dfrac{n}{d}) = 1$

\problem{1}Найти НОК и НОД двух чисел (без калькулятора и Очира!):

а) $2^7\cdot3^7\cdot5^8$ и $3^{73}\cdot7^{51}\cdot2^{64}\cdot11^5$

б) $777$ и $808$

в) $270, \  450 \ \text{и} \ 555.$

\problem{2} У Анджура Аркадьевича 615 чокопаев и 861 бабуфохов. Нужно составить из них одинаковые наборы,
причем так, чтобы раздать их наибольшему количеству детей. Сколько чокопаев и сколько бабуфохов должно быть тогда в наборе? Сколько детей получат подарки?

\problem{3} Приведите пример, когда равенство  $(a, b, c)[a, b, c] = abc$  не выполнено. 

\problem{4}а)Найдите НОД$(a^m - 1, \ b^m - 1 )$\\
б) Найдите  НОД$(111...111, 11...11)$  – в записи первого числа 100 единиц, в записи второго – $60.$

\problem{5}Проверьте является ли дробь $\dfrac{7x-33}{4x+25}$ несократимой. Ответ я не знаю, мне лень считать.

\problem{6}Сумма двух чисел равна 1070, а НОК равно 44844. Найти эти числа.

\problem{7} Алина посчитала НОК всех чисел от 1 до 1000, а Глеб — всех чисел от 501 до 1000. У кого получилось больше и во сколько раз?

\problem{8}Докажите равенство  $(F_n, \ F_m) = F(m, n)$.


\problem{9}\textit{Мини гробик от Очира.} Найдите НОД всех трёхзначных чисел составленных только из 1, 2, 3.

\problem{10}\textbf{МЕГА ГРОБ ОТ РАСУЛА.} Число вида $f_k = 2^{2^{k}} + 1 $ называют числом Ферма. Докажите, что разные числа Ферма всегда взаимно просты. 